%LTeX: language=fr-FR
\documentclass[12pt,a4paper]{article}

% ---------- Langue et encodage ----------
\usepackage[french]{babel}
%\usepackage[T1]{fontenc}
%\usepackage[utf8]{inputenc}

% ---------- Mise en page ----------
\usepackage[margin=2.2cm]{geometry}
\usepackage{array,booktabs,tabularx}
\usepackage{amsmath,amssymb}
\usepackage{xcolor}
\usepackage{tcolorbox}
\tcbuselibrary{skins,breakable}
\usepackage{enumitem}
\usepackage{graphicx}
\usepackage{tikz}
\usepackage{pgfplots}
\pgfplotsset{compat=1.18}
\usepackage{fancyhdr}
\usepackage{makecell}
\usepackage{pgffor}
\usepackage{currfile}
\usepackage{lastpage}

\graphicspath{ {../../Images/} }
\input{\currfiledir ../def.tex}
\def \Document {Activité 1}


% ---------- Couleurs ----------
\definecolor{bleuF}{RGB}{47,87,163}   
\definecolor{orH}{RGB}{242,170,60}    
\definecolor{cadre}{RGB}{47,87,163}

% ---------- Boîtes ----------
\newtcolorbox{objectif}{
  colback=bleuF!6, colframe=cadre, boxrule=0.6pt, arc=2pt,
  left=6pt, right=6pt, top=4pt, bottom=4pt,
  title={\textbf{Objectif}}, fonttitle=\bfseries,
  coltitle=cadre, colbacktitle=bleuF!12
}
\newtcolorbox{bilan}{
  colback=orH!8, colframe=cadre, boxrule=0.6pt, arc=2pt,
  left=6pt, right=6pt, top=4pt, bottom=4pt,
  title={\textbf{Bilan}}, fonttitle=\bfseries,
  coltitle=cadre, colbacktitle=orH!20,
  breakable
}
\newtcolorbox{correctionbox}{
  colback=green!5, colframe=green!40!black, boxrule=0.6pt, arc=2pt,
  left=6pt, right=6pt, top=4pt, bottom=4pt,
  breakable
}

% \checkmark symbol fallback (amssymb provides \checkmark)

% ---------- Ligne de réponse pointillée ----------
\newcommand{\reponse}[1]{%
  \par\vspace{2mm}
  \foreach \i in {1,...,#1}{\rule{0pt}{0pt}\dotfill\\[3mm]}%
}

\begin{document}


\pagestyle{fancy}

\fancyhf{} % clear existing header/footer entries
\fancyhead[R]{ {\small \Niveau}}
\fancyhead[L]{ {\small \ChapitreCourt}}
\fancyhead[C]{ {\small \Document}}
%\fancyhead[R]{\large \textbf{\textbf{\Document\ (B) -- 1h}}}
\fancyfoot[L]{ A. Fernandez }
\fancyfoot[R]{{\small \thepage\  / \pageref{LastPage}}}

% ---------- Repère vierge pour placer les points ----------
\newcommand{\repere}[2]{% #1 : légende de l'axe des ordonnées, #2 : ymax
\begin{center}
\begin{tikzpicture}
\begin{axis}[
  width=0.8\textwidth, height=7.5cm,
  xmin=-0.5, xmax=7.5, ymin=0, ymax=#2,
  xtick={0,...,7}, ytick distance=10, minor y tick num=1,
  grid=both, major grid style={gray!45}, minor grid style={gray!20},
  axis lines=left, xlabel={$n$}, ylabel={#1},
  ylabel style={font=\small}, xlabel style={at={(1,0)}, anchor=west},
  ]
\end{axis}
\end{tikzpicture}
\end{center}}

\setlength{\parindent}{0pt}
\newcommand{\trou}[1]{\underline{\hspace{#1}}}

\begin{objectif}
Découvrir la notion de suite : calculer des termes, passer d'un terme au suivant, représenter les termes dans un repère.
\end{objectif}

\section*{Partie A -- Le nombre d'enfants inscrits}

Une crèche municipale ouvre à la rentrée 2026 avec $20$ enfants inscrits. La mairie prévoit d'accueillir $6$ enfants de plus à chaque rentrée.

On note $u(n)$ le nombre d'enfants inscrits à la rentrée de l'année $2026+n$. Ainsi $u(0)=20$.

\begin{enumerate}[label=\textbf{\arabic*.}, leftmargin=1.3em]
\item Compléter le tableau suivant.
\begin{center}
\renewcommand{\arraystretch}{1.8}
\begin{tabular}{|c|*{7}{>{\centering\arraybackslash}p{1.2cm}|}}
\hline
Année & 2026 & 2027 & 2028 & & & & \\
\hline
$n$ & $0$ & $1$ & $2$ & & & & \\
\hline
$u(n)$ & $20$ & & & & & & \\
\hline
\end{tabular}
\end{center}

\item Que représente le nombre $u(4)$~?
\reponse{1}

\item Compléter : $u(n+1) = u(n) + \trou{1.5cm}$.

\item Sans compléter le tableau plus loin, calculer $u(10)$.\\
En déduire une formule : $u(n) = 20 + \trou{1.5cm} \times n$.
\reponse{2}

\item Placer les points de coordonnées $(n~;~u(n))$ pour $n$ allant de $0$ à $6$. A-t-on le droit de relier les points~? Pourquoi~?
\repere{Nombre d'enfants}{80}
\reponse{1}
\end{enumerate}

\section*{Partie B -- Le budget des repas}

En 2026, la crèche dépense $40\,000$~€ pour les repas. Ce budget augmente de $10\,\%$ chaque année.

On note $v(n)$ le budget des repas, en euros, pour l'année $2026+n$. Ainsi $v(0)=40\,000$.

\begin{enumerate}[label=\textbf{\arabic*.}, leftmargin=1.3em]
\item Un prix de $100$~€ augmente de $10\,\%$. Quel est le nouveau prix~? Par quel nombre a-t-on multiplié $100$~?

En déduire : augmenter une quantité de $10\,\%$, c'est la multiplier par \trou{1.5cm}.
\reponse{2}

\item Compléter le tableau suivant, puis compléter : $v(n+1) = \trou{1.5cm} \times v(n)$.
\begin{center}
\renewcommand{\arraystretch}{1.8}
\begin{tabular}{|c|*{5}{>{\centering\arraybackslash}p{1.6cm}|}}
\hline
$n$ & $0$ & $1$ & $2$ & $3$ & $4$ \\
\hline
$v(n)$ & $40\,000$ & & & & \\
\hline
\end{tabular}
\end{center}

\item Peut-on calculer $v(10)$ aussi rapidement que $u(10)$~? Pourquoi~?
\reponse{2}
\end{enumerate}

\begin{bilan}
\begin{itemize}[leftmargin=1.2em]
\item Une \textbf{suite} est une liste de nombres numérotés par les entiers $0$, $1$, $2$, \dots{} Le nombre $u(n)$ est le \textbf{terme de rang} \trou{1.2cm}.
\item \textbf{Formule explicite} : on calcule un terme directement à partir de son rang $n$. Exemple : $u(n) = \trou{3cm}$.
\item \textbf{Relation de récurrence} : on calcule un terme à partir du terme \trou{3cm}. Il faut connaître le premier terme. Exemple : $v(0) = \trou{2cm}$ et $v(n+1) = \trou{3cm}$.
\item On représente une suite par un \textbf{nuage de points} $(n~;~u(n))$. On ne \trou{2cm} pas les points.
\end{itemize}
\end{bilan}

\end{document}
